\documentclass[10pt]{article}
\usepackage{tmlr}
\usepackage{amsmath,amssymb,booktabs,graphicx,tabularx,longtable}
\usepackage[hidelinks]{hyperref}
\usepackage{microtype}
\usepackage{adjustbox,placeins}
\graphicspath{{figures/}}
\newcommand{\ind}{\mathbf{1}}
\title{Evaluating Deadline-Limited Communication in Multi-Agent Reinforcement Learning:\\Delivery, Computation, and Control}
\author{\name Anonymous authors}
\begin{document}
\maketitle

\begin{abstract}
Lower traffic at maintained return does not identify better message selection. We examine this attribution problem under a shared serial link with reception deadlines and cached messages. On matched Hallway policies, rotating access reduces attempted traffic by 80\% while preserving all 400 observed delivery, action and reward histories. Changing a feasible fresh sender subset changes subsequent return in 13.9--18.9\% of sampled states across four ten-seed checkpoint--condition combinations, yet an advantage over uniform selection is sensitive to state versus episode weighting. Shared initialization eliminates the observed slow-link return contrast despite later delivery and action differences. CPU preparation changes discrete packet opportunities, and broadcast probes reveal usable information that acting policies do not fully exploit. Finally, an optimal no-message Hallway controller qualifies the learned communication advantage. Together, these results show why traffic--return comparisons need delivery, initialization and computation controls before gains can be attributed to message selection.
\end{abstract}


\section{Introduction}
Communication-efficient multi-agent reinforcement learning seeks to retain cooperative performance while reducing message exchange. VBC and TMC pursue this objective through selective or temporally controlled communication~\citep{zhang2019efficient,zhang2020succinct}; SchedNet learns which agents access a shared medium~\citep{kim2019schedule}. These objectives motivate reporting traffic alongside task return. Interpreting the resulting trade-off, however, requires distinguishing reduced transmission from better information selection.

An action deadline makes that distinction concrete. A packet can be generated and charged without arriving in time for the action. Removing such attempts reduces traffic without changing what the controller receives. Conversely, two rules with the same admission cap may deliver different packet counts because their computations consume different portions of the reception window. Even at the same received count, changing a sender may change cache initialization, later actions and the trajectory on which subsequent messages are evaluated. A traffic--return pair cannot distinguish these mechanisms.

\paragraph{Central claim.}
Under deadline-limited communication, traffic and episodic return alone do not identify the control value of a message-selection rule. We test the links from attempted traffic to delivery, from delivery to policy response, and from policy response to return. We use ten-seed Hallway policies with a shared GRU architecture and matched-checkpoint interventions. Round-Robin (RR) is the simple access reference; Deadline-Aware Communication (DAC) combines a capacity cap with Q-gap/transmitter-age priority.

The experiments follow that attribution chain. Matched Full--RR histories separate traffic savings from information changes. Same-count sender interventions and common-prefix branches separate sender choice from initial cache filling. CPU repeats test whether nominally matched channels offer the same delivery opportunities. Broadcast probes test downstream use, and an optimal no-message Hallway controller tests task necessity.

\newpage
\vspace*{-2.5\baselineskip}\vspace*{-6pt}\vspace*{-7mm}
\enlargethispage{2mm}
\section{Related Work}
\paragraph{From communication reduction to sender selection.}
Learned channels, aggregation and targeted messages support cooperative policies~\citep{foerster2016learning,sukhbaatar2016learning,das2019tarmac}, while IC3Net learns when to communicate~\citep{singh2019learning}. VBC studies communication efficiency through variance-based control~\citep{zhang2019efficient}. TMC reports reduced exchange with maintained performance and robustness to transmission loss~\citep{zhang2020succinct}; ETCNet explicitly imposes a bandwidth constraint~\citep{hu2023event}. These are communication-efficiency objectives. Our Full--RR identity comparison asks the narrower attribution question: can the same traffic reduction arise without changing any received information?

SchedNet learns shared-medium access and compares against round-robin scheduling~\citep{kim2019schedule}; IMAC addresses efficient communication through an information bottleneck~\citep{wang2020imac}. Graph attention, teammate modeling and importance-aware exchange add richer selection mechanisms~\citep{liu2020game,yuan2022incentive,huang2024qmnet}, alongside model-based sparsity and asynchronous scheduling~\citep{han2023model,han2025sparse,srinivasan2025asynchronous}. These approaches motivate our fixed-packet-count subset and common-prefix comparisons: an episodic scheduler difference can combine sender choice with a changed information history. We evaluate that ambiguity on a shared GRU/QMIX backbone~\citep{rashid2018qmix}, rather than reproducing or reranking those algorithms. \citet{zhu2024survey} review the broader design space.

\paragraph{From fresh delivery to useful control.}
DACOM learns how long an agent should wait for messages~\citep{yuan2023dacom}. We examine the complementary cost of preparing those messages and selecting their senders before reception begins. Age-of-information work measures update freshness~\citep{kaul2012real,liu2022optimizing,gong2022hierarchical,alibotaiken2026freshness}; task-oriented control scheduling connects transmission decisions to control objectives~\citep{ayan2022semantics}. Freshness and information use therefore require separate observations. \citet{lowe2019pitfalls} distinguish signaling from listening and show why coarse communication metrics can mislead. Our broadcast probes and Hallway interventions extend this evaluation question to finite delivery opportunities: which information reaches the action head, whether behavior changes, and whether that change improves return. The cited studies supply the design objectives and interpretation questions; our experiments test them within the channel and policies specified below.

\section{Evaluation Contract and Access Rules}
\subsection{Policy, channel and cost boundary}
Agents $i\in\{0,\ldots,N-1\}$ share a GRU with local input $[o_i^t;\operatorname{onehot}(a_i^{t-1});e_i]$, where $e_i$ is the agent identifier. A linear head maps hidden state $h_i^t$ to message $m_i^t$. Received messages condition the action head, not the recurrent update. Each message has eight float32 entries: payload $P=32$ B and modeled header 16 B, hence wire size $S=48$ B. Evaluation return is undiscounted $\sum_t r_t$; learning uses discount 0.99.

CPU inference uses one PyTorch thread and \texttt{perf\_counter\_ns}. Measured computation $C_t$ includes encoding, gate selection, serialization, cache operations, transport arithmetic, aggregation, action output and residual bookkeeping; environment stepping and logging are outside this boundary. Tensor-to-byte serialization is charged. If $P_{\mathrm{ready},t}$ is elapsed CPU time at packet readiness, a fixed reserve $R_{\rm post}=0.45$ ms defines $D_{\rm recv}=D_{\max}-R_{\rm post}$. The reserve does not replace actual post-reception computation.

Attempted senders $\mathcal S_t$ are queued by $(i-t)\bmod N$, regardless of priority score. For $K_t=|\mathcal S_t|$, bandwidth $B$ and fixed latency $\ell=0.05$ ms,
\begin{align}
 A_{k,t}&=P_{\mathrm{ready},t}+kS/B+\ell,\qquad d_{k,t}=\ind[A_{k,t}\le D_{\rm recv}],\\
 W_t&=\begin{cases}\min\{K_tS/B+\ell,[D_{\rm recv}-P_{\mathrm{ready},t}]_+\},&K_t>0,\\0,&K_t=0,\end{cases}\\
 T_{\rm step,t}&=C_t+W_t.\label{eq:step}
\end{align}
Accepted broadcasts reach every receiver except the sender. Every attempted packet is charged in full, including cutoff drops; each step flushes the queue. The model has no losses, collisions, retransmissions, acknowledgments or inter-step backlog. Computation and modeled waiting add without overlap: this is hybrid accounting, not measured radio latency. A deadline miss occurs when $T_{\rm step,t}>D_{\max}$; the cutoff cannot eliminate arbitrary CPU jitter.

\subsection{Cache and accounting}
Delivery replaces sender $j$'s cache with $m_j^t$ and its source timestamp with $t$; otherwise both persist. Initially caches are zero and invalid. Validity $v_j^t$ distinguishes never-received entries from fresh zero payloads; receiver age $\tau_j^t$ is time since the valid source timestamp. For receiver $i$,
\begin{equation}
 g_i^t=\frac{\sum_{j\ne i}v_j^t\lambda^{\tau_j^t}c_j^t}{\max(1,\sum_{j\ne i}v_j^t)}.
\end{equation}
The action head uses $[h_i^t;g_i^t;\log(1+\bar\tau_i^t);n_i^t/(N-1)]$, with valid-neighbor count $n_i^t$ and mean valid age $\bar\tau_i^t$. The denominator counts valid senders, not decay weights. Thus delivery changes content, aggregate scaling and metadata together; it is not a pure content intervention.

For measured steps $\mathcal T$ and received count $M_t$, offered load, payload goodput and waste are
\begin{equation}
 L_{\rm off}=\frac{S\sum_tK_t}{|\mathcal T|},\qquad G_{\rm pay}=\frac{P\sum_tM_t}{|\mathcal T|},\qquad
 \rho_{\rm drop}=1-\frac{\sum_tM_t}{\sum_tK_t}.
\end{equation}
Each broadcast is counted once, not per receiver. Waste is undefined when no packets are attempted. Return, load, goodput, cache validity/age, computation and deadline misses measure different outcomes.

\subsection{Admission and selection}
At elapsed admission time $E_t$, the cap combines link capacity and traffic credit:
\begin{align}
 K_{\rm link}(E_t)&=\left\lfloor [D_{\max}-R_{\rm post}-E_t-\ell]_+B/S\right\rfloor,\label{eq:cap}\\
 c_t^-&=\min(c_{\max},c_{t-1}+b),\quad K_{\max,t}=\min\{N,K_{\rm link}(E_t),\lfloor c_t^-/S\rfloor\},\\
 c_t&=c_t^- - SK_t.
\end{align}
Here $b=120$ B/step, $c_{\max}=144$ B and initial credit is zero. This bounds cumulative attempted bytes, not each step's bytes. Candidate packing is charged even if a post-packing cap recheck rejects it. Later bookkeeping can still invalidate an earlier capacity estimate.

Full attempts all senders; Periodic selects $(t+i)\bmod2=0$. RR uses the cap and advances its cursor by the admitted count. DAC uses the same cap but ranks batched local scores
\begin{equation}
 u_i^t=(Q_{i,(1)}^{\rm local}-Q_{i,(2)}^{\rm local})+0.1a_i^{\rm tx,t}.
\end{equation}
The local head has zero communication features; the gap is zero with only one legal action. The threshold is zero and ties use agent ID. Transmitter age starts at zero, resets on an attempt (even if dropped), and increments otherwise. It differs from receiver age. RR/DAC use $\lambda=0.9$; Full/Periodic use $\lambda=1$. Identity controls explicitly match decay. DAC requires all agents' scores at a central scheduler: CPU cost is measured, but physical score-collection traffic is not. Its score is a heuristic, not a measure of marginal receiver utility.

\section{Experimental Questions and Protocol}
\subsection{Questions and cohorts}
The comparisons separate four potential explanations of a traffic--return result: admission, sender identity and initialization, selector computation, and the acting policy's use of information. Independently trained Hallway policies supply descriptive comparisons (Table~\ref{tab:primary}); interventions use frozen primary RR/DAC weights (Tables~\ref{tab:new_fresh}--\ref{tab:new_prefix}). The broadcast probe uses a separate cohort (Table~\ref{tab:broadcast_use}). These estimates are not pooled. Training and environment details are in Appendix~\ref{app:settings}.

\newpage
\vspace*{-2.5\baselineskip}\vspace*{-6pt}\vspace*{-7mm}\vspace*{8.8pt}
\subsection{Fresh-weight Hallway campaign}
Five agents form groups of sizes two and three. A successful group contributes 10 and both groups can yield a maximum episodic return of 20; the finite horizon is 14 steps. For each of ten seeds, we initialize six configurations from fresh weights: no communication, ideal fresh communication, Full under the modeled channel, Periodic, capacity-capped Round-Robin (RR), and DAC. Each configuration trains for 1,500 updates in each of Slow/5 ms, Medium/2 ms and Medium/5 ms, yielding 180 jobs. Learning uses the Hallway Oracle/BC support specified in Appendix~\ref{app:settings}. This is not demonstration-free learning. Equal update counts do not imply equal environment transition counts, because policies can terminate episodes at different times.

Final checkpoints are evaluated on 50 paired held-out resets in Slow/5 ms, Medium/2 ms, Medium/5 ms and Medium/1.5 ms: 720 seed-level cells and 36,000 episode returns. All seeds remain in primary summaries. There is no best-checkpoint or successful-seed selection. Ideal Full overrides delivery failure to supply fresh messages; it is a reference condition rather than an optimal-return bound. RR uses the same credit, link cap and cache decay as DAC, advances its cursor by admitted sender count, and avoids local-Q priority computation. Equal caps do not imply identical actual bytes when CPU costs or trajectories differ.

\subsection{Statistical and timing controls}
Means and sample SD are computed across independent training seeds. Paired bootstrap resamples seed differences 10,000 times. Episode repetitions do not increase the number of independent training seeds. A supplemental post-hoc exact two-sided sign-flip analysis uses Holm correction separately over three matched-condition Ideal--No-Comm contrasts and all twelve DAC--RR train/deployment contrasts. Sign-flip inference requires exchangeability under the null, or symmetric paired differences; it is not a distribution-free guarantee. Unadjusted bootstrap intervals and corrected p-values are labeled separately.

CPU evaluation uses one thread, excludes two warmup episodes, disables automatic garbage collection during rollouts and collects outside them. The post-reception reserve is fixed at 0.45 ms. No CPU timing claim is derived from GPU optimization. A separate matched-checkpoint diagnostic fixes preparation time at 0.5 ms for admission and transport while recording actual computation separately. This isolates a common reception opportunity; its virtual timing is not deployable latency.

\subsection{Additional hidden-source tasks}
Broadcast and XOR are custom N=5 fresh-bit diagnostic tasks, not ports of independent standard benchmarks. At every step one source observes a fresh bit (broadcast), or two sources each observe one fresh bit (XOR). Followers predict the bit or parity. Fresh independent bits and absence of source information in follower observations imply expected no-communication follower accuracy of one half, regardless of recurrent memory; an ideal bit-sharing oracle can attain one. This does not guarantee that the learned ideal-link policy solves the task.

The registered campaign trains five methods (None, Ideal, Periodic, RR, DAC), ten seeds and two tasks for 1,000 updates each: 100 jobs, no demonstrations and BC weight zero. Training is at Medium/2 ms; final CPU evaluation uses Medium/5 ms, Medium/2 ms and Slow/5 ms with 50 paired resets. The horizon is 20; accuracy equals return divided by 20. Before interpreting scheduler effects, we check whether at least eight of ten ideal-link policies reach 75\% accuracy. The learning gate prevents an unsuccessful ideal-link learner from being used to diagnose communication utility.

\clearpage
\vspace*{-2.5\baselineskip}\vspace*{-6pt}\vspace*{-7mm}
\section{Results}
\vspace{-6pt}
\begin{table}[!ht]\centering\setlength{\tabcolsep}{3pt}\small
\caption{Primary Hallway results: independently trained configurations, matched training/deployment condition, all ten seeds, mean and sample SD. Learned rows are Oracle-assisted, 1,500 updates per job, with measured CPU readiness and 50 evaluation resets per seed. Reference rows have no training or messages: random actions use ten action-RNG seeds on the same 50 primary resets (8200000--8200049), with SD across RNG means, not training seeds. The local-clock value is verified over all 1,024 initial states. Neither reference uses the link, so its value is repeated across columns.}\label{tab:primary}
\vspace{2.5mm}
\begin{adjustbox}{max width=\linewidth}\begin{tabular}{lrrr}\toprule
Configuration & Slow/5 ms & Medium/2 ms & Medium/5 ms\\\midrule
No-Comm & $3.96\pm2.27$ & $4.00\pm2.38$ & $4.26\pm1.98$\\
Ideal Full & $11.44\pm2.98$ & $12.20\pm3.04$ & $11.54\pm3.22$\\
Full & $6.60\pm3.15$ & $7.70\pm1.81$ & $11.54\pm3.26$\\
Periodic & $6.54\pm3.53$ & $5.66\pm1.99$ & $7.88\pm4.16$\\
RR & $6.12\pm1.69$ & $5.64\pm1.99$ & $7.88\pm2.25$\\
DAC & $4.72\pm3.13$ & $6.38\pm0.82$ & $7.04\pm3.52$\\
\midrule Random actions (reference) & $0.22\pm0.20$ & $0.22\pm0.20$ & $0.22\pm0.20$\\
Constructive local clock (reference) & 20 & 20 & 20\\
\bottomrule\end{tabular}\end{adjustbox}\end{table}
\begin{table}[!ht]\centering\setlength{\tabcolsep}{3pt}\small
\caption{Paired seed effects. CIs are unadjusted percentile bootstrap intervals. Supplemental post-hoc two-sided sign-flip p-values use Holm correction over three utility contrasts or all twelve train/deployment DAC--RR contrasts.}\label{tab:paired}
\vspace{3mm}
\begin{adjustbox}{max width=\linewidth}\begin{tabular}{llrrr}\toprule
Condition & Contrast & Difference & 95\% CI & Holm p\\\midrule
Medium/2 ms & Ideal--No-Comm & +8.20 & [+5.48,+10.64] & 0.0059\\
Medium/5 ms & Ideal--No-Comm & +7.28 & [+4.88,+9.38] & 0.0059\\
Slow/5 ms & Ideal--No-Comm & +7.48 & [+4.92,+9.94] & 0.0059\\
Medium/2 ms & DAC--RR & +0.74 & [-0.96,+1.96] & 1.0000\\
Medium/5 ms & DAC--RR & -0.84 & [-3.02,+1.26] & 1.0000\\
Slow/5 ms & DAC--RR & -1.40 & [-3.06,+0.08] & 1.0000\\
\bottomrule\end{tabular}\end{adjustbox}\end{table}
\FloatBarrier
\subsection{Communication utility versus sender-selection advantage}
Comparisons with and without a communication channel establish a learned-policy performance gap, as used in communication learning~\citep{foerster2016learning,sukhbaatar2016learning}. Following the finer distinction raised by \citet{lowe2019pitfalls}, we first measure that gap and then test what it can identify. Table~\ref{tab:primary} compares independently trained policies and two untrained reference policies. Ideal Full improves on No-Comm by 7.48, 8.20 and 7.28 return points in Slow/5 ms, Medium/2 ms and Medium/5 ms, respectively. All corresponding unadjusted paired intervals exclude zero; the supplemental Holm-adjusted sign-flip p-values are 0.00586 (Table~\ref{tab:paired}). These results describe communication utility under privileged Oracle teaching. The exhaustive Oracle audit in Appendix~\ref{app:settings} identifies conflicting teacher labels for identical local histories; therefore the comparison also depends on the information requirements of supervision. It does not isolate a same-policy message-content effect.

In contrast, DAC minus RR is $-1.40$, $+0.74$ and $-0.84$ in the same matched conditions, with every paired interval including zero. The intervals leave both positive and negative selection effects plausible. Among cross-condition deployments, Medium/5 ms training followed by Slow/5 ms deployment gives DAC $+2.60$, unadjusted CI $[0.98,4.34]$; its corrected p-value over all twelve selector comparisons is 0.1875. The favorable deployment cell does not survive that multiplicity control.

\subsection{Admission efficiency does not identify the best controller}
The communication-reduction objective of VBC and TMC~\citep{zhang2019efficient,zhang2020succinct} admits two explanations: retaining useful information with fewer attempts, or selecting information that improves control. Our cutoff model lets us separate them. Unrestricted Full attains mean returns 6.60, 7.70 and 11.54 in Slow/5 ms, Medium/2 ms and Medium/5 ms. The corresponding RR means are 6.12, 5.64 and 7.88, and DAC means are 4.72, 6.38 and 7.04. These point estimates do not show a return improvement from admission control. Admission must therefore be judged on both attempted traffic and controller performance. The model charges entire attempted packets and flushes the queue each step; attempted-byte savings are not a measured radio-energy saving.


For the matched Slow/5 ms trained policies, Full offers 240 wire B/step for 32 payload B/step and 80\% waste. Periodic offers 120.14 B/step for the same payload goodput and approximately 60.05\% waste. Both RR and DAC offer 48 B/step, retain 32 payload B/step, and exhibit zero observed drops, but return is 6.12 for RR and 4.72 for DAC. Capacity capping explains the shared communication saving; it cannot uniquely explain DAC's sender selection.

This reduction follows the modeled capacity constraint. At preparation time $P_{\mathrm{ready},t}$, the link permits at most
\begin{equation}
 K_t^{\mathrm{feasible}}=\left\lfloor \frac{B_{\rm link}[D_{\rm recv}-P_{\mathrm{ready},t}-\ell]_+}{S}\right\rfloor
\end{equation}
arrivals of equal-sized packets. Avoiding attempts beyond this bound prevents modeled late packets if preparation accounting remains exact. The saving follows directly from finite service capacity. Later CPU work and timing variability can still invalidate an earlier cap estimate: matched Medium/2 ms DAC has approximately 0.031\% waste. There it offers 92.23 B/step versus RR's 95.46 B/step; the comparison is same-cap, not exactly equal offered load.

\subsection{Different traffic can produce identical control histories}
Figure~\ref{fig:seed_distribution} displays individual seed means, making dispersion and overlap visible beyond the aggregate table.
Full and RR deliver identical sender sets at all 6,238 common observed Slow/5 ms steps, but their original cache decay differs ($\lambda=1$ versus $0.9$) and their independently trained weights differ. The return gap is therefore not an identified training-noise floor. A separate fixed-readiness identity check uses ten slow-trained Full checkpoints, 20 paired resets each, and matched decay at both $\lambda=1$ and 0.9. All 400 paired episodes reproduce exactly the delivered sender histories, cache tensors, actions, rewards and return under Full versus RR. Full still offers five times the traffic. In these single-packet conditions, fewer attempts preserve every observed control outcome: admission improves traffic efficiency without improving sender choice. Periodic and RR share a Medium/5 ms mean of 7.88 but have different individual seed returns and zero identical delivered sets among 5,412 common steps. Equal means do not establish duplicate policies.

\begin{table}[!t]\centering\setlength{\tabcolsep}{3pt}\footnotesize
\caption{Exploratory sensitivity of DAC--RR seed-level inference. TOST margin is $\pm1$ return, selected after observing earlier results; MDE is an approximate 80\% power sensitivity based on observed paired SD. No confirmatory equivalence claim follows.}\label{tab:sensitivity}
\vspace{3mm}
\begin{adjustbox}{max width=\linewidth}\begin{tabular}{lrrrrr}\toprule
Condition & Difference & t 95\% CI & BCa 95\% CI & TOST p & MDE 80\\\midrule
Slow/5 ms & -1.40 & [-3.29,0.49] & [-3.26,-0.08] & 0.678 & 2.59\\
Medium/2 ms & +0.74 & [-1.11,2.59] & [-1.77,1.78] & 0.379 & 2.54\\
Medium/5 ms & -0.84 & [-3.47,1.79] & [-2.96,1.36] & 0.447 & 3.60\\
\bottomrule\end{tabular}\end{adjustbox}\end{table}
Table~\ref{tab:sensitivity} shows that BCa and t intervals disagree on zero exclusion for Slow/5 ms, demonstrating method sensitivity at ten seeds. No matched condition passes the exploratory $\pm1$ equivalence criterion. The intervals are unadjusted and the sensitivity analysis does not replace the multiplicity-controlled primary analysis.
\begin{figure}[!t]\centering\includegraphics[width=\linewidth]{seed_distribution.pdf}
\caption{Ten independently trained seed means for each of Full, Periodic, RR and DAC in each matched condition; horizontal marks show the mean. Episode repetitions are not independent training replicates.}
\label{fig:seed_distribution}
\end{figure}

\section{Where Traffic--Return Attribution Breaks Down}
SchedNet's comparison with round-robin access makes sender selection an explicit design question~\citep{kim2019schedule}. Figure~\ref{fig:controlled} separates that question from initial information acquisition and from the computation that determines delivery opportunities.
\subsection{Sender choice and initial cache filling have different effects}
To test selection at the same information budget, we use the ten primary training seeds separately for RR-trained and DAC-trained weights. At prespecified even timesteps along both policies' trajectories, we enumerate every feasible one- or two-sender fresh-message subset, keeping the pre-decision state, policy weights and packet count fixed. Each branch resumes the original access rule thereafter. The best--worst suffix-return range and best-subset regret are hindsight diagnostics, not evidence of performance achievable by an online scheduler. All 2,654 no-op and 2,654 forced-original-subset checks passed. Both compare the full recorded suffix trace: actions, rewards, termination flags, environment states, delivered sender sets, source timestamps and accounted traffic.

The shared-prefix control transmits with RR until every sender is valid, then clones the state and switches to RR or DAC. All prefix traffic and modeled time are included. Episodes ending before full coverage remain in total-return summaries with zero suffix contrast; survivor-only effects are separately reported. This comparison tests whether the episodic scheduler difference survives shared initial information. Cached-message methods such as TMC~\citep{zhang2020succinct} motivate separating initial acquisition from later refresh.


Eligible senders are all five agent IDs, including inactive agents; this diagnostic does not filter them by activity or score. The cloned snapshot is immediately before encoding and gate selection, and $k(x)$ is the factual number of delivered packets at that state (one or two here). Define $\mathcal A_{k(x)}=\{A\subseteq\{0,\ldots,N-1\}:|A|=k(x)\}$, the set of all sender subsets of cardinality $k(x)$. All $\binom{5}{k(x)}$ subsets are enumerated under fixed readiness. Forced selection is applied at both admission calls. Transmitter age resets for the actually forced attempted senders, credit is charged for those attempts, and the RR cursor advances by their count, not by their identities. The original access rule resumes from that changed scheduler state; later history changes are part of the total effect.
\begin{align}
 V_x(A)&=R_{t:T}(\text{deliver }A\text{ now; resume original access rule}),\\
 \mathrm{Range}(x)&=\max_{A\in\mathcal A_{k(x)}}V_x(A)-\min_{A\in\mathcal A_{k(x)}}V_x(A),\\
 \mathrm{Regret}(x)&=\max_{A\in\mathcal A_{k(x)}}V_x(A)-V_x(A_{\rm factual}).
\end{align}
Effect is $100\,\mathbf1[\mathrm{Range}(x)>0]$. Within each training seed we average equally over its sampled states, without an intermediate episode average; the reported mean then weights the ten seeds equally. Longer factual episodes therefore contribute more states within a seed. Every interval in Table~\ref{tab:new_fresh} is the unadjusted descriptive Student-$t$ interval of these ten seed means (9 degrees of freedom), not an interval treating branches as independent observations. Regret uses a hindsight best subset and is not achievable online performance. Appendix~\ref{app:robustness} directly compares factual and uniform-subset outcomes at the same received count. Candidate branch counts exclude no-op, forced-original and withholding checks; per-seed denominators accompany the supplementary data.

\begin{table}[!t]\centering\setlength{\tabcolsep}{3pt}\footnotesize
\caption{Feasible-subset effects on primary checkpoints. Each row uses ten frozen primary training seeds, ten evaluation resets per seed and fixed 0.5 ms readiness; Effect is the seed-mean percentage of sampled states with nonzero best--worst return range, with an unadjusted 95\% t interval.}\label{tab:new_fresh}
\vspace{3mm}
\begin{adjustbox}{max width=\linewidth}\begin{tabular}{lllll}\toprule
Condition & Weights & States & Subsets & Effect (\%) [CI]\\\midrule
Slow/5 ms & RR & 646 & 3,230 & 18.86 [15.33, 22.40]\\
Slow/5 ms & DAC & 687 & 3,435 & 13.94 [9.45, 18.44]\\
Medium/2 ms & RR & 661 & 6,610 & 16.97 [9.99, 23.94]\\
Medium/2 ms & DAC & 660 & 6,600 & 18.16 [11.13, 25.19]\\
\bottomrule\end{tabular}\end{adjustbox}\end{table}
\begin{table}[!t]\centering\setlength{\tabcolsep}{3pt}\footnotesize
\caption{Shared-prefix effects on frozen primary checkpoints: ten training seeds, 20 resets per seed and fixed 0.5 ms readiness. Both differences are DAC minus RR under the indicated weights; cold starts use empty caches, while prefix branches share an RR initialization. Early termination is retained; intervals use ten seed means.}\label{tab:new_prefix}
\begin{adjustbox}{max width=\linewidth}\begin{tabular}{lllll}\toprule
Condition & Weights & Cold diff. & Prefix diff. [95\% CI] & Survivors\\\midrule
Slow/5 ms & RR & -0.850 & +0.000 [0.000,0.000] & 183/200\\
Slow/5 ms & DAC & +0.200 & +0.000 [0.000,0.000] & 190/200\\
Medium/2 ms & RR & -1.700 & +0.350 [-0.235,0.935] & 199/200\\
Medium/2 ms & DAC & -0.200 & -0.350 [-1.025,0.325] & 200/200\\
\bottomrule\end{tabular}\end{adjustbox}\end{table}

\begin{figure}[!t]\centering\includegraphics[width=\linewidth]{diagnostics.pdf}
\caption{Primary-checkpoint diagnostics. Left: each point is one seed's mean feasible-subset suffix-return range. Middle: Slow/5 ms only, cold and common-prefix DAC--RR effects on the same frozen policy, with early termination retained. Right: medium link, primary Medium/2 ms-trained RR/DAC checkpoints, mean delivered packets over ten seeds under measured (solid) or fixed (dashed) readiness. Seed-specific curves are retained in the supplementary summary; this panel shows means without uncertainty bands. These are distinct interventions, not pooled training cohorts.}\label{fig:controlled}
\end{figure}
\FloatBarrier
\paragraph{Weighting and remaining-task sensitivity.}
Replaying all 400 factual trajectories verifies the 2,654 sampled states and finds no simultaneous-arrival penalty in those factual trajectories. Outcome-resolved waiting occupies 476 of 5,290 factual steps. Excluding resolved sampled states raises effect rates to 16.48--25.92\%; equal-episode instead of equal-state weighting gives 14.76--21.79\% (Table~\ref{tab:robust_strata}). These rates remain conditional on the sampled checkpoints and histories. With equal-state weighting, the directly paired factual-minus-uniform contrast is positive with an unadjusted t interval excluding zero for Slow/DAC and Medium/RR, but not the other two cells. With equal-episode weighting, all four intervals include zero. Thus a positive selection advantage over uniform is not robust to this change of estimand, even though nonzero subset effects remain. Consequential sender choices therefore exist, but the current priority rule does not robustly exploit them across weighting choices. Appendix~\ref{app:robustness} gives the aggregation and counterfactual sensitivity boundary.

\paragraph{Initialization versus later refresh.}
After the common RR prefix, suffix return changes in 0 of 400 Slow/5 ms episodes and 15 of 400 Medium/2 ms episodes. Early termination is retained (Table~\ref{tab:new_prefix}). In the slow condition, 126 pairs nevertheless differ in delivered sets and 94 in actions on common suffix steps. Of 373 survivors, 38 have no unresolved group, 296 have one and 39 have two, with nine steps remaining. An executed global-state action reference verifies unused attainable reward in 335/400 pairs. Thus identical messages, universal termination and reward saturation do not explain the zero return contrast. The same finite-horizon return can thus hide different later information and actions, even when additional reward remains attainable.

\subsection{Selector computation changes the available channel}
DACOM treats the benefit of waiting for messages as a control decision~\citep{yuan2023dacom}. Here, sender-selection computation consumes time before waiting begins, changing the service opportunity itself. At Medium/2.25 ms, three 48-byte packets require readiness within 310 microseconds, after a 450-microsecond reserve, 50-microsecond propagation and 1.44 ms serialization. RR/DAC mean gate times are 128.7/208.4 microseconds; mean deliveries are 2.175/1.999 packets under measured readiness and 2.000/2.000 under fixed 500-microsecond readiness. Readiness tails and bookkeeping, not mean gate cost alone, determine service opportunities. Thus a nominally identical link does not present both implementations with identical reception opportunities.

\paragraph{Independent-process timing check.}
Three independent sequential processes on the same pinned CPU repeat all ten primary seeds at 2.0, 2.25 and 2.5 ms, using three evaluation resets per cell rather than the original ten. At 2.25 ms, RR delivers 2.213, 2.005 and 2.003 packets/step, while DAC delivers 2.000 in all three processes. The original packet-count advantage is therefore not a stable effect magnitude. On identical frozen-policy/input snapshots, DAC's additional gate cost averages 98.5, 99.4 and 100.7 microseconds across these processes, respectively. The extra computation persists, while its conversion into packet loss depends on timing tails and available credit.

\subsection{Delivered information need not be fully exploited}
\citet{lowe2019pitfalls} separate informative signaling from behavioral use. Table~\ref{tab:broadcast_use} applies that distinction after deadline-constrained delivery: we compare an unchanged acting policy with a probe of the aggregate it actually receives. Under source-aware selection, held-out decoding of the actual receiver aggregate reaches 99.91--100\%, while the unchanged policies achieve 64.17--78.44\%. Thus a lack of accessible aggregate information cannot be the sole explanation in these cells. The bottleneck includes how accessible information is used, rather than source delivery alone. The learning check passes for broadcast (10/10 ideal-link seeds reach 75\% accuracy) but fails for XOR (0/10). XOR contrasts are excluded from scheduler conclusions. Appendix~\ref{app:probe} specifies the separate held-out probe protocol.

\begin{table}[htbp]\centering\footnotesize
\caption{Broadcast delivery versus use: fixed 0.5 ms readiness, ten-seed mean accuracy (\%) on the same held-out ten episodes. Normal and source-aware acting policies share weights; the probe is fitted separately on ten disjoint episodes. Source-aware selection preserves the factual delivered count. Paired probe gains and protocol are in Appendix~\ref{app:probe}.}\label{tab:broadcast_use}
\vspace{3mm}
\begin{adjustbox}{max width=\linewidth}\begin{tabular}{lllll}\toprule
Condition & Weights & Normal policy & Source-aware policy & Source-aware aggregate probe\\\midrule
Slow/5 ms & RR & 51.90 & 78.44 & 100.00\\
Slow/5 ms & DAC & 50.27 & 70.94 & 100.00\\
Medium/2 ms & RR & 54.98 & 64.17 & 99.91\\
Medium/2 ms & DAC & 53.35 & 65.09 & 100.00\\
\bottomrule\end{tabular}\end{adjustbox}\end{table}

\subsection{A learned communication gap does not establish task necessity}
The channel-versus-no-channel comparison used in learned communication~\citep{foerster2016learning,sukhbaatar2016learning} measures the outcome of a training procedure. To test the stronger interpretation that the task requires messages, we construct a controller using only local distance and a public clock. It attains return 20 with zero communication on all 1,024 Hallway starts. This establishes solvability under the public-clock and local-information assumptions detailed in Appendix~\ref{app:local_clock}, not learning success of the neural No-Comm baseline. Learned Ideal--No-Comm gaps therefore concern the specified training protocol rather than task-level communication necessity.


\section{Discussion: Evaluating Useful Communication}
\paragraph{Separate an engineering saving from an explanation of that saving.}
The 80\% Full--RR reduction is a real attempted-traffic saving under the specified accounting rule. Its mechanism is removing packets that cannot affect the action, rather than choosing better received information. Communication-reduction objectives~\citep{zhang2019efficient,zhang2020succinct} remain useful, but traffic and return alone do not distinguish these explanations. A delivered-history identity control identifies when admission achieves the saving without changing control.

\paragraph{Compare selection after matching the opportunity to select.}
Round-robin scheduling is already a relevant reference in SchedNet~\citep{kim2019schedule}. Our results show why it must be matched at more than the nominal cap: realized packet counts, cache initialization and computation can differ. Once packet count is fixed, some sender choices change return, yet factual-over-uniform effects depend on the weighting of sampled states. After common slow-link initialization, subsequent delivery and action differences produce no observed return change. Selection quality should therefore be evaluated at matched states and packet counts, followed by episodic evaluation that includes initialization and actual cost.

\paragraph{Measure what the policy does with delivered information.}
Broadcast provides a concrete instance of the signaling--listening distinction~\citep{lowe2019pitfalls}: a receiver aggregate supports near-perfect decoding while the acting policy remains substantially less accurate. Conversely, Hallway supports optimal control without messages despite a learned Ideal--No-Comm gap. Together, these results separate information availability, its learned use, and task necessity. They explain why neither a decodable message nor a channel-on return gain is sufficient evidence that a scheduler selects the most useful information.

\paragraph{Scope of the evidence.}
The primary policies use Oracle-assisted five-agent Hallway; broadcast and XOR are custom tasks, and XOR fails the ideal-link learning criterion. A separate SMACv2 pilot also fails its preset criterion in Terran and Protoss (three Ideal Full seeds each, approximately 100,000 environment steps; final win rates 3.125--9.375\%); no scheduler comparisons followed. Thus standard-benchmark replication remains unresolved. Equal training updates need not mean equal environment steps. The serial channel excludes losses, collisions, retries, backlog and compute/transport overlap. CPU repeats use one host, and physical centralized score-collection traffic is uncharged. Fixed-state branches share training seeds, whereas prefix interventions measure total history changes rather than pure coverage mediation. The experiments diagnose attribution in these implementations; they do not rerank the cited native algorithms or establish a deployed-network guarantee.

\section{Conclusion}
Traffic reduction, useful delivery and improved control are different results. Under the tested deadline model, rotating access removes 80\% of attempts while preserving 400 Full-communication trajectories. Same-count interventions show that sender identity can matter, but the tested priority rule does not establish a general selection advantage. Shared initialization, CPU preparation and downstream information use each change how a traffic--return comparison should be read. Evaluations should therefore trace attempted messages through timely delivery and policy response to return, using matched histories for traffic savings and matched packet opportunities for selection quality. This turns a communication-efficiency comparison into evidence about its mechanism.

\clearpage\bibliographystyle{tmlr}\bibliography{references}
\clearpage\appendix
\section{Essential Reproducibility Details}\label{app:settings}
\paragraph{Primary Hallway learning.} Adam learning rate is $5\times10^{-4}$, discount 0.99, global gradient clip 10, and hard target synchronization every 100 updates. A fresh policy has shared GRU width 64, message width 8 and QMIX mixing; Double-Q targets use the online network to select the next action and the target network to evaluate that action. The online replay holds 2,000 episodes; 50 Oracle episodes populate a separate demonstration pool. Each 16-episode batch contains 4 demonstration and 12 online episodes, with replacement only while a stratum is too small. One online episode is collected per update; 1,500 updates therefore need not represent equal environment-step counts. The original worker uses exploration $1-0.95u/1499$ at zero-based update $u$. The TD objective adds 0.5 times demonstration cross entropy. Target synchronization and exploration count updates, not environment transitions. All final checkpoints are evaluated, without successful-seed filtering.

\paragraph{Hallway observations and teacher.}\label{app:local_clock} Each agent observes its own integer distance and active flag. The encoder additionally receives one-hot agent ID and previous action; the mixer state concatenates all distance/active pairs. Distances initialize independently and uniformly in $\{1,2,3,4\}$; group IDs are $[0,0,1,1,1]$. Actions wait, reduce distance by one, or increase distance by one up to four. A group whose members all reach zero together earns 10 and becomes inactive; an asynchronous partial arrival fails that group. If both groups arrive simultaneously, the step first earns +20 and then subtracts 30, giving net reward -10. An episode terminates naturally only when both groups have succeeded or both have failed; mixed success/failure reaches the 14-step horizon. Exact simultaneous-arrival semantics are given in the \emph{Environment aliasing and snapshot isolation} paragraph.

The privileged teacher moves agents toward distance one, then admits the lowest-index ready group. Across all 1,024 starts it obtains return 20; 39 of 120 distinct agent-specific local observation and previous-action history keys map to conflicting teacher actions. This is a count of distinct keys, not a BC error rate or a demonstrated cause of the learned return gap. A different no-message controller moves to distance one and uses public group-specific arrival steps 3 and 4. Given a synchronized reset clock, known roster and maximum initial distance, it obtains return 20 on all 1,024 starts (5,120 steps). This is exhaustive solvability, not neural learning success. BC applies to all agents at filled demonstration timesteps, including inactive agents, but not padding or online episodes.

\paragraph{Environment aliasing and snapshot isolation.}
The local variable \texttt{last\_state} aliases the environment's mutable distance array; it is not a saved pre-action copy. During simultaneous two-group arrival, the nominal rollback therefore retains the post-move zero distances. Both group-active flags remain false, agent-active flags are restored true, and the win counter returns to zero. No further group reward is evaluated, and the episode reaches the horizon. This is a preserved task-specific transition rule, not an alias between stored transitions or independent branches. Direct checks verify the net -10 reward and subsequent timeout. Adapter observations and global states are new arrays; episode construction copies them into arrays, replay insertion/sampling deep-copies records, and branch cloning deep-copies mutable environment/harness state while sharing immutable model parameters. Added tests mutate later environment states, replay samples and branch caches and verify the stored/factual data remain unchanged.

\paragraph{Replay consistency.} Replay stores full episodes beginning with an empty cache and zero initial recurrent state, not arbitrary chunks. Nonzero initial caches or source timestamps are rejected. The online and target networks each unroll their own recurrent states and regenerate their own messages under current respective parameters. Both use the same recorded delivery events. Consequently, the cache retained between events belongs to that network's latent space; it is not a stale stored vector encoded by older parameters. Next-state bootstrap uses the recorded next-observation delivery event, not a freshly sampled gate. Termination and horizon truncation both suppress bootstrap. This is off-policy evaluation conditional on replayed access history; it is not a claim that current gates would recreate the recorded behavior.

\paragraph{Measurement and artifacts.}
Slow and medium links are 20,000 and 100,000 B/s (decimal units). The host is an AMD Ryzen Threadripper 3960X, Python 3.13.12, PyTorch 2.11.0+cu128 and NumPy 2.4.4. Sequential repeat processes use the same allowed CPU, one inference thread, two discarded warmups and automatic GC disabled during rollouts; seed/selector orders are counterbalanced. Identical-input cost checks clone policy, recurrent/cache and scheduler state outside the timer. Fixed readiness replaces preparation time for admission/transport only, while CPU work is separately recorded. These are not extra training seeds or cross-machine replications.

The companion \texttt{supplementary.zip} provides source, configurations, diagnostic logs, seed aggregates and provenance. Its scope manifest lists included checkpoints and omitted intermediate arrays; commands regenerate the receiver features. Full-suffix identity tests compare actions, rewards, termination, delivery, cache timestamps and traffic. ``Registered'' refers to configurations recorded before their runs, not public preregistration.

\section{Weighting and Counterfactual Checks}\label{app:robustness}
Define the paired factual-minus-uniform contrast at state $x$ as
\[
\Delta_x=V_x(A_{\rm factual})-|\mathcal A_{k(x)}|^{-1}\sum_{A\in\mathcal A_{k(x)}}V_x(A).
\]
This enumerates a one-decision uniform reference, not a trained random scheduler. Equal-state weighting averages states within each seed; equal-episode weighting first averages states within episodes, then episodes within seeds. Both weight ten seeds equally; intervals use unadjusted Student-t inference with 9 degrees of freedom. Empty strata are not zero-imputed. Table~\ref{tab:robust_strata} shows the change of estimand explicitly. The 224 sampled states with no unresolved group have zero range. No factual episode incurred the simultaneous-arrival penalty. One of 19,875 saved counterfactual suffix totals is negative; nonnegative totals can still conceal intermediate penalties, so these checks do not remove all dependence on Hallway's transition rules.

\begin{table}[htbp]\centering\setlength{\tabcolsep}{3pt}\footnotesize
\caption{Fresh-subset sensitivity. Effect columns are percentages: original equal-state weighting, equal-episode weighting, and equal-state weighting restricted to unresolved tasks. All three weight seeds equally. Direct factual minus uniform-subset suffix return uses the identical sampled state and received count; unadjusted 95\% t intervals over ten seeds. Episode weighting first averages sampled states within an episode.}\label{tab:robust_strata}
\begin{adjustbox}{max width=\linewidth}\begin{tabular}{lllllll}\toprule
Condition & Weights & State & Episode & Unresolved & Paired state [CI] & Paired episode [CI]\\\midrule
Slow/5 ms & RR & 18.86 & 21.79 & 24.19 & -0.056 [-0.277, 0.166] & -0.178 [-0.526, 0.171]\\
Slow/5 ms & DAC & 13.94 & 14.76 & 16.48 & 0.182 [0.008, 0.356] & 0.130 [-0.064, 0.324]\\
Medium/2 ms & RR & 16.97 & 19.14 & 25.92 & 0.235 [0.100, 0.371] & 0.195 [-0.027, 0.418]\\
Medium/2 ms & DAC & 18.16 & 19.48 & 22.21 & -0.145 [-0.439, 0.149] & -0.249 [-0.636, 0.139]\\
\bottomrule\end{tabular}\end{adjustbox}\end{table}

\paragraph{Remaining opportunities.}
All 800 common prefixes were replayed. A privileged action reference schedules active groups at distinct feasible arrival times; its predicted reward matches 3,406 executed suffix steps. Slow/5 ms has 65, 296 and 39 endpoints with attainable remaining rewards 0, 10 and 20, respectively; 335 admit more reward than either learned suffix. Medium/2 ms has 1, 260 and 139 such endpoints and 361 with unused opportunity. This action-changing reference tests saturation, not message-only causality.

\section{Broadcast Learning and Receiver Probes}\label{app:probe}
Broadcast fixes source ID 0; XOR uses IDs 0 and 1. Sources observe independent fresh signed bits, source flags and indices; followers receive no current bit, source action or reward as an observation. All agents see normalized timestep. Reward is the fraction of followers correctly predicting the bit or parity, over horizon 20. The mixer sees bits and time, while decentralized heads see only local input and messages. All 100 jobs use Adam $5\times10^{-4}$, 1,000 updates, batch 16, replay 1,000, hidden/message sizes 64/8 and no Oracle/BC. The original learning check uses 50 resets (9500000--9500049) and measured readiness; broadcast passes in 10/10 seeds and XOR in 0/10.

The probe instead fixes readiness at 0.5 ms and uses 20 resets (14200000--14200019). Source-aware access forces source 0 into a subset with the factual delivered count at each cloned state. Source identity is supplied; this is not learned generic selection. The subsequent histories may differ, and equal delivered count does not guarantee equal offered traffic. On all 20 resets, slow-link RR/DAC normal source delivery is 20.0/17.0\%, increasing to 100\% with source-aware access. Diagnostic selection CPU overhead is unmeasured.

Episodes 14200000--14200009 fit per-receiver ridge probes; 14200010--14200019 evaluate them and the acting-policy scores in Table~\ref{tab:broadcast_use}. Normalization uses fit episodes only, coefficient $10^{-3}$ and an unpenalized intercept. A pre-hook records the actual 74-dimensional utility-head input; the aggregate is independently reconstructed with validity, decay and self-exclusion. The 40,000 replayed steps preserve source bits, delivery, count and reward; a separate 6,000-step replay checks wire-byte accounting.

Source-aware aggregate-probe minus acting-policy gains are 21.56 [7.19,35.94] and 29.06 [11.90,46.23] percentage points for Slow/5 ms RR and DAC, and 35.74 [23.91,47.57] and 34.91 [19.34,50.49] for Medium/2 ms RR and DAC (unadjusted 95\% t intervals of ten seed means). These probes use training labels unavailable to the acting head. Decodability rules out absent accessible aggregate information as the sole bottleneck in these cells; it does not identify how to make the policy exploit it.
\end{document}
